\begin{table}[!htbp]
\centering\small
\caption{Notación utilizada en la demostración.}
\label{tab:proof-notation}
\begin{tabular}{@{}p{.29\linewidth}>{\raggedright\arraybackslash}p{.65\linewidth}@{}}
\toprule
Símbolo & Significado \\
\midrule
$n,k$ & Orden del tablero inicial y número de filas vacías en el tablero residual. \\
$M,K$ & Corte $M=\lceil\delta n\rceil$ entre los dos tramos y tamaño terminal $K$. \\
$B_t,B^{(k)},B_\tau$ & Tablero tras $t$ inserciones, tablero con $k$ filas vacías
($B^{(k)}=B_{n-k}$) y tablero en la parada. \\
$\mathcal F_t,\mathcal F^{(k)}$ & Historia revelada, indexada por inserciones o por filas restantes,
con $\mathcal F^{(k)}=\mathcal F_{n-k}$. \\
\addlinespace[.22em]
$Q(B),L(B),g(B)$ & Conteo de completaciones, logaritmo del predictor y diferencia
$g(B)=\log Q(B)-L(B)$ cuando $Q(B)>0$. \\
$p_e,\ell_f,h_f$ & Peso de Bethe de la casilla $e$, masa de la diagonal $f$ y su vacancia
$h_f=1-\ell_f$. \\
\addlinespace[.22em]
$\mu,\nu,P^{\rm B}$ & Ley uniforme de revelado, ley adaptativa basada en el predictor
y ley auxiliar de Bethe. \\
$\tau,\rho(B,i),A_\tau$ & Tiempo de la primera salida o de la llegada al tamaño $K$, lo que ocurra
primero, cociente de normalización de la fila y logaritmo acumulado del cociente
$A_\tau=\sum_{t<\tau}\log\rho(B_t,i_t)$. \\
\addlinespace[.22em]
$F,G$ & Llegada al tamaño $K$ con regularidad hasta ese tablero inclusive, y subsuceso
$G=F\cap\{Q(B^{(K)})>0\}$. \\
$D_K$ & Suceso en el que dos cruces vacíos distintos en el tamaño $K$ comparten una diagonal. \\
\bottomrule
\end{tabular}
\end{table}
